\documentclass[a4,10pt]{aleph-notas}

% -- Paquetes adicionales
\usepackage{enumitem}
\usepackage{aleph-comandos}

% -- Datos
% El siguiente texto presenta la definición de la jerarquía acumulativa de von Neumann
% y demuestra que sus niveles cubren todos los conjuntos (la clase universal).
% Adaptación y exposición didáctica basada en el tratamiento estándar (ver Corolario 1.3.6).
\institucion{Proyecto Alephsub0}
\asignatura{Matemáticas de la Información}
\tema{Jerarquía de von Neumann: definición y universalidad}
\autor{Andrés Merino}
\fecha{Abril 2026}

% Fuente tipográfica y logotipo
\fuente{montserrat}
\logouno[4.5cm]{Logos/LogoAlephsub0-02}
\definecolor{colordef}{cmyk}{0.81,0.62,0.00,0.22}

% -- Comandos adicionales
\newtheorem*{prob}{Problema}
    \tcolorboxenvironment{prob}{%
        color=colordef,recuadrost,colback=colordef!7,drop fuzzy shadow
    }


\begin{document}


\encabezado

\noindent
El siguiente problema corresponde al Ejemplo 1.3.7 del libro Introduction to Cardinal Aritmetic de M. Holz. Primero, presentemos algunas definiciones:

\begin{defi}
    Definimos, de manera recursiva:
    \begin{enumerate}
        \item $V_0=\varnothing$.
        \item $V_{\alpha+1}=\mathcal{P}(V_\alpha)$ para todo ordinal $\alpha$.
        \item Si $\gamma$ es un ordinal límite, entonces $V_\gamma=\bigcup\{V_\beta:\beta<\gamma\}$.
    \end{enumerate}
    Los conjuntos $V_\alpha$ se llaman niveles de la jerarquía de von Neumann y la función $(V_\alpha:\alpha\in\mathrm{ON})$ se denomina jerarquía (o jerarquía acumulativa).
\end{defi}

\begin{lem}
    Sean \(\alpha\) y \(\beta\) ordinales tales que \(\alpha\leq \beta\). Entonces
    \[
        V_\alpha\subseteq V_\beta.
    \]
\end{lem}

\begin{prob}
    Probar que la clase 
    \[
        \bigcup\{V_\alpha:\alpha\in\mathrm{ON}\}
    \]
    es la clase universal.
\end{prob}

\begin{proof}
Por reducción al absurdo, supongamos que existe un conjunto \(x\) tal que
\[
    x\notin V := \bigcup\{V_\alpha:\alpha\in\mathrm{ON}\}.
\]
De esta manera, se tiene que la clase
\[
    V^{c}=\left\{x:x\notin V\right\} \neq \varnothing.
\]
Por el axioma de regularidad, existe \(a\in V^c\) tal que
\[
    a\cap V^c=\varnothing.
\]
Luego, si \(y\in a\), entonces \(y\notin V^c\). Por lo tanto,
\[
    y\in V.
\]
Así, para cada \(y\in a\), existe un ordinal \(\alpha_y\) tal que
\[
    y\in V_{\alpha_y}.
\]

Por el axioma de reemplazo, la clase
\[
    \{\alpha_y:y\in a\}
\]
es un conjunto de ordinales. Entonces existe un ordinal \(\beta\) tal que
\[
    \alpha_y<\beta
\]
para todo \(y\in a\). De aquí se sigue que
\[
    y\in V_\beta
\]
para todo \(y\in a\). Por tanto,
\[
    a\subseteq V_\beta.
\]
Luego,
\[
    a\in \mathcal{P}(V_\beta)=V_{\beta+1}.
\]
En particular, \(a\in V\), lo cual contradice que \(a\in V^c\).

Por lo tanto, \(V^c=\varnothing\). En consecuencia,
\[
    V=\bigcup\{V_\alpha:\alpha\in\mathrm{ON}\}
\]
es la clase universal.
\end{proof}

\begin{obs}
    En la demostración anterior no es necesario utilizar el axioma de elección. En efecto, para cada \(y\in a\), como \(y\in V\), existe al menos un ordinal \(\alpha\) tal que
    \[
        y\in V_\alpha.
    \]
    Como la clase de los ordinales está bien ordenada por \(\in\), existe un único ordinal mínimo con esta propiedad. Es decir, podemos definir
    \[
        \alpha_y=\min\{\alpha\in\mathrm{ON}:y\in V_\alpha\}.
    \]
    De esta manera, la asignación \(y\mapsto \alpha_y\) queda definida, sin realizar elecciones arbitrarias.
\end{obs}

\end{document}